% =============================================================================
%  Part III: From Single Scatterer to Effective Medium  (Sections 17--18)
% =============================================================================

\part*{Part III: From Single Scatterer to Effective Medium}

\section{Homogenisation and Effective Medium Theory}

\subsection{The Space-Filling Requirement}

The ultimate goal of computing single-scatterer T-matrices is to build
effective medium theories for heterogeneous elastic materials.  The idea is
simple: model the heterogeneous medium as a collection of scatterers, each
characterised by its T-matrix, and then average over the ensemble to obtain
effective material properties.

The critical requirement for such a homogenisation procedure is that the
scatterers must \emph{tile} the medium completely.  If they do not, the
material between scatterers is unmodelled, leading to systematic
errors \citep{Willis1977,Hashin1983}.

Cubes are the simplest three-dimensional shape that tiles $\mathbb{R}^3$
without gaps or overlaps.  This was famously proven to be optimal for convex
shapes by the Keller--Minkowski theorem, and the practical importance for
effective medium theory was recognised early on by Hashin and
Shtrikman \citep{HashinShtrikman1962a,HashinShtrikman1962b,HashinShtrikman1963},
who used cubic partitions as the foundation of their variational bounds on
effective elastic moduli.

\subsection{From T-Matrix to Effective Moduli}

Given the cubic T-matrix~\eqref{eq:Tcube}, we can construct the effective
elastic moduli of a composite material as follows.  Consider a background
medium partitioned into $N$ non-overlapping cubes, each with (possibly
different) material properties.  In the long-wavelength limit, the
effective stiffness tensor satisfies \citep{ZellerDederichs1973,Willis1977}:
\begin{equation}
  C_{ijkl}^{\text{eff}} = C_{ijkl}^{(0)}
  + \frac{1}{V}\sum_{n=1}^{N}
    \Delta c_{ijuv}^{(n)}\,A_{uv,kl}^{(n)},
\label{eq:Ceff}
\end{equation}
where $C^{(0)}$ is the background stiffness, $\Delta c^{(n)}$ is the
contrast of the $n$-th cube, and $A^{(n)}$ is the amplification operator
derived from its T-matrix.  In the dilute limit (non-interacting
scatterers), each $A^{(n)}$ is computed independently as in Section~15.

\subsection{Self-Consistent Effective Medium Schemes}

More sophisticated approximations account for scatterer-scatterer
interactions:

\textbf{Coherent Potential Approximation (CPA):}  Each cube is embedded in
the (unknown) effective medium rather than the background, leading to a
self-consistent equation for $C^{\text{eff}}$ \citep{ZellerDederichs1973,Gubernatis1977}.
The T-matrix formulation is particularly natural here: the CPA condition
is simply that the average T-matrix vanishes, $\langle T\rangle = 0$.

\textbf{Differential Effective Medium (DEM):}  Scatterers are added
incrementally to the background, with the effective medium updated at each
step \citep{Norris1985}.  The cubic T-matrix provides the ``update kernel''
for each infinitesimal addition.

\textbf{Kramers--Kronig consistency:}  Because the T-matrix is derived
from the full self-consistent Lippmann--Schwinger equation (not the Born
approximation), the resulting effective moduli are causal and satisfy
Kramers--Kronig relations \citep{Willis1981}, making them suitable for
broadband modelling.

\subsection{Cubic Symmetry of the Effective Medium}

Even when the individual cubes have isotropic material properties
($\Delta\lambda$, $\Delta\mu$, $\Delta\rho$), the effective medium inherits
\emph{cubic symmetry} from the scatterer geometry.  This is a direct
consequence of $T_3\neq 0$ in~\eqref{eq:Tcube}: the effective shear
modulus splits into $\mu^{*,\text{diag}}$ and $\mu^{*,\text{off}}$, giving
the effective stiffness tensor three independent elastic constants (bulk
modulus, and two shear moduli) rather than the two of an isotropic medium.

The cubic anisotropy parameter~\eqref{eq:cubicaniso} quantifies this
splitting.  At low frequencies it is $O(\omega a)^7$---very small---but it
grows rapidly with frequency, becoming significant as $ka$ approaches unity.

\subsection{Connection to Hashin--Shtrikman Bounds}

The classical Hashin--Shtrikman variational
bounds \citep{HashinShtrikman1962a,HashinShtrikman1962b,HashinShtrikman1963}
on effective elastic moduli are derived using a comparison medium and a
polarisation stress field.  The bounds are achieved by specific
microstructural geometries: the Hashin assemblage (concentric spheres)
achieves the bulk modulus bound, but the shear modulus bound requires more
complex geometries.

The cubic T-matrix approach provides a complementary, wave-based route to
effective properties that naturally includes frequency dependence and
finite-wavelength corrections absent from static variational bounds.  In
the zero-frequency limit, the cubic effective moduli computed here should
lie within the Hashin--Shtrikman bounds, providing a consistency check.


\section{Born Verification (Cube)}

In the Born limit $a\to 0$:
\begin{equation}
  \lim_{a\to 0}\Delta\rho^{*,c} = \Delta\rho, \qquad
  \lim_{a\to 0}\Delta\lambda^{*,c} = \Delta\lambda, \qquad
  \lim_{a\to 0}\Delta\mu^{*,\text{diag}} = \Delta\mu, \qquad
  \lim_{a\to 0}\Delta\mu^{*,\text{off}} = \Delta\mu.
\label{eq:Bornlim_cube}
\end{equation}
The cubic anisotropy vanishes:
$\lim_{a\to 0}(\Delta\mu^{*,\text{diag}} - \Delta\mu^{*,\text{off}}) = 0$.
All limits verified analytically with SymPy.
